En un modelo idealizado de una zona de almacenamiento, la cantidad de material \(m\) se trata como una variable continua, en kilogramos. Entra material a un caudal constante \(u\ge0\), que puede elegirse. El equipo de extracción retira material al caudal
\[
q(m)=\frac{\mu m}{K+m},\qquad
\mu=12\,\mathrm{kg/min},\quad K=6\,\mathrm{kg}.
\]
La extracción se anula cuando no queda material y se aproxima a un caudal máximo cuando aumenta la cantidad almacenada. No hay otras entradas ni salidas. Inicialmente hay \(4\,\mathrm{kg}\), y se quiere evitar que el almacenamiento supere \(9\,\mathrm{kg}\).

Construya el modelo y determine todos los caudales de entrada que permiten respetar ese límite para cualquier \(t\ge0\). Justifique la elección mediante el signo de la tasa de cambio, sin obtener una solución explícita. Explique qué ocurre con la cantidad almacenada cuando se utiliza el mayor caudal admisible y por qué cualquier caudal mayor hace superar el límite en un tiempo finito.
